\section{Background}\label{sec:background}

\subsection{Before RAG: The Limits of Standalone Language Models}
Historically, applying large language models to specialized tasks relied heavily on the knowledge stored directly within the model's own weights during training. While this approach—often enhanced through supervised fine-tuning or reinforcement learning—allows models to reason and format answers well, it presents significant drawbacks for knowledge-intensive fields like medicine. A standalone model cannot easily update its knowledge without expensive retraining. More importantly, when asked about highly specific or rare medical conditions, models relying solely on internal memory are prone to hallucination, generating answers that sound confident but lack factual accuracy.

\subsection{RAG and the Challenge of Domain-Specific Search}
To address the limitations of standalone models, Retrieval-Augmented Generation (RAG)~\cite{lewis2021rag} was developed. RAG improves a model's accuracy by searching an external database for relevant facts and providing that text to the model before it answers. While fundamental RAG pipelines effectively retrieve general knowledge, applying them to highly specialized, massive databases presents unique challenges. Standard RAG relies on matching a user's raw question to target documents in a shared space. However, in medical contexts, user queries are often brief or phrased differently than formal clinical texts, creating a gap in meaning. As a result, standard search models frequently struggle to connect simple queries to documents containing rare medical terminology. This vocabulary mismatch leads to poor search results, leaving the reasoning agent without the necessary context to form an accurate answer.

\subsection{Hypothetical Document Embeddings (HyDE)}
To overcome the vocabulary limitations of standard RAG, advanced systems implement Hypothetical Document Embeddings (HyDE) as a search extension~\cite{gao2022hyde}. HyDE bypasses the gap between the user's simple query and the complex source document by shifting the approach to document-to-document matching. Instead of directly searching with the user's initial prompt, the system uses a language model to generate a hypothetical document in response to the question. While this generated text might contain factual inaccuracies, it successfully captures the expected language patterns, clinical vocabulary, and keywords naturally associated with the medical topic. Searching the database with this medically styled document significantly improves the retrieval of relevant information compared to using the raw question.

\subsection{The Medline Dataset}
In this project, the external knowledge base is grounded in Medline, a premier bibliographic database containing millions of references to journal articles in the life sciences and biomedical fields\cite{medlineplus}. Medline serves as an authoritative source of medical knowledge, covering everything from common treatments to rare diseases and complex pharmacology. However, the language used in Medline articles is highly technical, dense, and formal. This makes it an ideal testing ground for advanced retrieval techniques like HyDE, as successfully navigating Medline requires bridging the gap between everyday user questions and deeply specialized medical literature.
